\vspace{-6mm}
\section{Related Work}
\label{sec:related-work}

\subsection{Sensorless 3D US Reconstruction}
\label{sec:related-work:sensorless-3d-us-reconstruction}

Tracked 3D US reconstruction uses calibrated optical or electromagnetic
tracking to localize B-mode frames before interpolation or compounding
\cite{huang2017}.
Sensorless reconstruction instead infers inter-frame geometry from image
sequences through learned geometry-localization models, removing the need for
external pose measurements.
Within this paradigm, Prevost et al.~\cite{prevostDeepLearningSensorless2017}
pioneered pairwise probe-motion regression from neighboring frames, while
Guo et al.~\cite{guoSensorlessFreehand3D2020} extended it with contextual
information from multiple consecutive frames.
Subsequent methods incorporated anatomical shape priors,
contextual-contrastive learning, online adaptation, and long-range sequence
modeling to improve trajectory estimation
\cite{luoSelfContextShape2021,guoUltrasoundVolumeReconstruction2023,
luoRecONOnlineLearning2023,li2024f}.
Despite these advances, most approaches still recover frame-wise or inter-frame
6-DoF motion without explicitly exploiting regular sweep structure, retaining
unnecessary geometric complexity
\cite{huangMMSEReconstruction2008,liTUSREC2024}.
In contrast, our AFGM exploits such regularity to avoid dense pose recovery,
providing a compact and globally consistent spatial scaffold for robust
downstream 3D reconstruction.

\vspace{-3mm}
\subsection{Non-Gaussian Methods for 3D US Reconstruction}
\label{sec:related-work:non-gaussian-methods-for-3d-us-reconstruction}

Initially, 3D US volumes were commonly reconstructed from localized B-mode
slices through interpolation or compounding
\cite{mozaffariFreehand3DUltrasound2017}.
With the rise of Neural Radiance Fields (NeRF)
\cite{mildenhallNeRFRepresentingScenes2020}, implicit neural representations
became a prominent alternative to discrete volumetric reconstruction.
Song et al.~\cite{song2022a} introduced an MLP-based implicit representation
for freehand carotid reconstruction.
ImplicitVol~\cite{yeungSensorlessVolumetricReconstruction2024} jointly
optimized a continuous volume and sensorless slice locations, while
RapidVol~\cite{eidRapidvolRapidReconstruction2025} factorized the representation
to accelerate optimization.
For ultrasound-specific neural rendering, Ultra-NeRF
\cite{wysockiUltraNeRFNeuralRadiance2024} incorporated acoustic physics into
NeRF; NeRF-US~\cite{dagliNeRFUSRemovingUltrasound2024} further introduced
geometric priors to suppress artifacts in uncontrolled acquisitions, while
Dou et al.~\cite{douContinuousBernoulliDistribution2024} explored probabilistic
backscatter modeling.
Although flexible, repeated neural-field evaluations increase optimization and
rendering cost, motivating efficient explicit representations.

\vspace{-3mm}
\subsection{Gaussian Methods for 3D US Reconstruction}
\label{sec:related-work:gaussian-methods-for-3d-us-reconstruction}

Compared with NeRF, 3D Gaussian Splatting (3DGS) replaces neural-field queries
with anisotropic Gaussians and differentiable splatting, enabling
efficient, high-fidelity reconstruction and rendering
\cite{kerbl3DGaussianSplatting2023}.
Building on this backbone, natural-image methods adapt 3DGS to challenging
acquisition settings through task-specific modules and optimization while
retaining its splatting formulation.
For motion-blurred inputs, BAD-Gaussians jointly models blur and camera motion,
while DiET-GS combines event constraints with diffusion priors
\cite{zhaoBADGaussiansBundle2024,leeDiETGSDiffusionPrior2025}.
For sparse views, FSGS introduces proximity-guided unpooling, while LoopSparseGS
combines initialization, depth alignment, and sparse-friendly sampling
\cite{zhuFSGSRealTimeFewShot2025,baoLoopSparseGSLoopBased2025}.
For dynamic scenes, Ex4DGS uses static-dynamic separation and temporal
interpolation, SWinGS adopts sliding windows, and USplat4D models per-Gaussian
uncertainty
\cite{leeFullyExplicitDynamic2024,shawSWinGSSlidingWindows2024,
guoUncertaintyMattersDynamic2026}.

\begin{figure*}[t]
    \centering
    \includegraphics[scale=0.5]{figures/pdf/methods_overview.pdf}
    \vspace{-3mm}
    \caption{
        Overview of ULGSplat.
        (a) Acquisition-constrained Frame Geometry Model (AFGM) constructs canonical geometry and provides $\mathcal{T}^{\mathrm A}$ to (b) Gaussian Scene Representation maintaining primitives $\mathcal{G}$.
        (c) Physics-aware Local-Global Renderer (PLGR) uses $\mathcal{G}$ and $\mathcal{T}^{\mathrm A}$ for finite-slice local echo formation and path-dependent global transmission, rendering B-mode images whose GT residuals optimize $\mathcal{G}$.
        (d) Contribution-guided Gaussian Adaptation Strategy (CGAS) periodically updates $\mathcal{G}$ via residual-guided densification and path-aware pruning.
    }
    \vspace{-4mm}
    \label{fig:methods-pipeline}
\end{figure*}


When observations depart from pinhole-camera images, the imaging process
changes, making the original 3DGS renderer and adaptation no longer directly
applicable.
For slice-based volumetric imaging, FruitNinja reconstructs object interiors
from cross-sectional supervision, while GaussianPile introduces slice-aware
Gaussian placement and finite-thickness projection
\cite{wuFruitNinja3DObject2025,kongGaussianPileUnifiedSparse2026}.
Compared with generic slice imaging, ultrasound additionally involves acoustic
propagation, attenuation, and reflection, raising the need for specialized
Gaussian designs.
Accordingly, UltraGauss replaces perspective projection with probe-plane
intersection and in-plane aggregation
\cite{eidUltraGaussUltrafastGaussian2025}; UltraGS combines depth-aware
Gaussians with physics-inspired acoustic rendering
\cite{yangUltraGSRealTimePhysicallyDecoupled2025}; and UltraG-Ray further models
attenuation and reflection through Gaussian ray casting
\cite{duelmerUltraGRayPhysicsBasedGaussian2026}.
However, existing methods emphasize either slice-local echo formation
or path-dependent propagation, while Gaussian adaptation remains incompatible
with these ultrasound-specific rendering roles.
We propose PLGR to jointly model local echo and global transmission,
together with CGAS to adapt the Gaussian population from their rendering
contributions, enabling efficient, high-fidelity B-mode synthesis with
ultrasound-consistent representation control.
